% Example: Multi-line aligned equations with \begin{align}
% Demonstrates how the align environment typesets several equations
% with alignment at a chosen character (typically &), and how
% align* suppresses equation numbers.
\documentclass[12pt, a4paper]{article}
\usepackage{amsmath, amssymb, amsthm}
\usepackage{fontspec}
\begin{document}

\section*{Aligned Multi-line Equations}

The \texttt{align} environment aligns equations at the \texttt{\&}
marker. Each line is numbered:

\begin{align}
  f(x) &= (x + 2)(x - 3) \\
       &= x^2 - 3x + 2x - 6 \\
       &= x^2 - x - 6
\end{align}

Use \texttt{align*} to omit numbers:

\begin{align*}
  (a + b)^2 &= (a + b)(a + b) \\
            &= a^2 + ab + ba + b^2 \\
            &= a^2 + 2ab + b^2
\end{align*}

Multiple independent aligned blocks can share the environment:

\begin{align}
  x + y &= 5 \label{eq:sum} \\
  x - y &= 1 \label{eq:diff}
\end{align}

From \eqref{eq:sum} and \eqref{eq:diff} we obtain \( x = 3 \) and
\( y = 2 \).

\end{document}
